% Motor: XeLaTeX, fijado en .latexmkrc y en la receta del editor.
\documentclass[aspectratio=169,10pt,c]{beamer}

% !TEX root = slides.tex
% Preambulo comun: estilo UTP, colores, fuentes, plantillas de beamer y
% comandos de apoyo didactico. Se carga desde slides.tex tras \documentclass.

% La opción "c" de beamer reparte el espacio libre 1:1.5 (arriba:abajo), por lo
% que el contenido queda por encima del centro. Igualamos ambos rellenos.
\makeatletter
\define@key{beamerframe}{c}[true]{%
  \beamer@frametopskip=0pt plus 1fill\relax%
  \beamer@framebottomskip=0pt plus 1fill\relax%
  \beamer@frametopskipautobreak=0pt plus .5\paperheight\relax%
  \beamer@framebottomskipautobreak=0pt plus .5\paperheight\relax%
}
\makeatother

\usepackage{graphicx}
\usepackage{tikz}
\usepackage{xcolor}
\usepackage{iftex}

\ifPDFTeX
  \usepackage[utf8]{inputenc}
  \usepackage[T1]{fontenc}
  \usepackage{lmodern}          % tipografia escalable para el respaldo con pdflatex
  \renewcommand{\familydefault}{\sfdefault}
\else
  \usepackage{fontspec}
  \IfFontExistsTF{Arial}
    {\setsansfont{Arial}[BoldFont={Arial Bold}]}
    {\setsansfont{Helvetica}[BoldFont={Helvetica Bold}]}
\fi

\usepackage[spanish,es-nodecimaldot]{babel}
\usepackage{microtype}
\usepackage{booktabs}
\usepackage{amssymb}
\usepackage{array}
\usetikzlibrary{arrows.meta,positioning,fit,calc,shapes.geometric}

\graphicspath{{assets/}}

% ---------- Identidad visual UTP ----------
\definecolor{utpRojo}{HTML}{C42B35}
\definecolor{utpNegro}{HTML}{111111}
\definecolor{utpGris}{HTML}{5F5F5F}
\definecolor{utpGrisMedio}{HTML}{9A9A9A}
\definecolor{utpGrisClaro}{HTML}{F1F1F1}
\definecolor{utpNaranja}{HTML}{C1670B}

\usetheme{default}
\setbeamercolor{normal text}{fg=utpNegro,bg=white}
\setbeamercolor{structure}{fg=utpRojo}
\setbeamercolor{frametitle}{fg=utpNegro,bg=white}
\setbeamercolor{title}{fg=utpNegro}
\setbeamercolor{subtitle}{fg=utpNegro}
\setbeamercolor{author}{fg=utpNegro}
\setbeamercolor{date}{fg=utpNegro}
\setbeamercolor{block title}{fg=white,bg=utpNegro}
\setbeamercolor{block body}{fg=utpNegro,bg=utpGrisClaro}
\setbeamercolor{alertblock title}{fg=white,bg=utpRojo}
\setbeamercolor{alertblock body}{fg=utpNegro,bg=utpRojo!7}
\setbeamercolor{exampleblock title}{fg=white,bg=utpGris}
\setbeamercolor{exampleblock body}{fg=utpNegro,bg=utpGris!8}

\setbeamerfont{title}{series=\bfseries,size=\fontsize{19.5}{22.5}\selectfont}
\setbeamerfont{subtitle}{size=\fontsize{11.5}{13.5}\selectfont}
\setbeamerfont{date}{series=\bfseries,size=\fontsize{11.5}{13.5}\selectfont}
\setbeamerfont{author}{size=\fontsize{10.5}{12.5}\selectfont}
\setbeamerfont{frametitle}{series=\bfseries,size=\large}
\setbeamerfont{block title}{series=\bfseries,size=\small}
\setbeamerfont{block body}{size=\normalsize}

\setbeamertemplate{navigation symbols}{}
\setbeamertemplate{headline}{}
\setbeamertemplate{blocks}[default]
\setbeamertemplate{itemize item}{\color{utpRojo}\small\raise.15ex\hbox{$\blacktriangleright$}}
\setbeamertemplate{itemize subitem}{\color{utpGris}\scriptsize$\bullet$}
\setbeamertemplate{enumerate item}{\color{utpRojo}\insertenumlabel.}
\setbeamertemplate{background canvas}{\color{white}\rule{\paperwidth}{\paperheight}}

% Título de diapositiva: negrita negra a la izquierda, logo UTP a la derecha
\setbeamertemplate{frametitle}{%
  \begin{tikzpicture}[remember picture,overlay]
    \node[anchor=north east,inner sep=0pt]
      at ([xshift=-7mm,yshift=-5.5mm]current page.north east)
      {\includegraphics[width=17mm]{utp-logo.png}};
  \end{tikzpicture}%
  \vspace{1.2mm}%
  \begin{minipage}[t]{.80\textwidth}\raggedright\insertframetitle\end{minipage}\par
  \vspace{1.1mm}{\color{utpRojo}\rule{\textwidth}{1.1pt}}
}

% Solo numeración discreta; la plantilla UTP no lleva pie de página
\setbeamertemplate{footline}{%
  \hfill{\color{utpGrisMedio}\scriptsize\insertframenumber}\hspace{6mm}\vspace{2.5mm}%
}

% Portada UTP
\setbeamertemplate{title page}{%
  \begin{tikzpicture}[remember picture,overlay]
    \node[anchor=center,inner sep=0pt] at (current page.center) {%
      \includegraphics[width=\paperwidth,height=\paperheight]{utp-cover-bg.jpeg}%
    };
    \node[anchor=north,align=center,text=utpNegro,inner sep=0pt,
      text width=145mm,font=\usebeamerfont{title}]
      at ([yshift=-25.2mm]current page.north) {\inserttitle};
    \node[anchor=north,align=center,text=utpNegro,inner sep=0pt,
      text width=142mm,font=\usebeamerfont{subtitle}]
      at ([yshift=-38.2mm]current page.north) {\insertsubtitle};
    \node[anchor=north,align=center,text=utpNegro,inner sep=0pt,
      font=\usebeamerfont{date}]
      at ([yshift=-50.4mm]current page.north) {\insertdate};
    \node[anchor=north,align=center,text=utpNegro,inner sep=0pt,
      font=\usebeamerfont{author}]
      at ([yshift=-58.4mm]current page.north) {\insertauthor};
  \end{tikzpicture}%
}

% ---------- Comandos de apoyo didáctico ----------
\newcommand{\concepto}[2]{%
  \begin{block}{#1}\large #2\end{block}}
\newcommand{\pregunta}[1]{%
  \begin{alertblock}{Decide y explica}#1\end{alertblock}}
\newcommand{\analogia}[1]{%
  \begin{exampleblock}{En palabras simples}\small #1\end{exampleblock}}
\newcommand{\enlarealidad}[1]{%
  \begin{block}{Cómo se ve en la realidad}\small #1\end{block}}
\newcommand{\ojo}[1]{%
  \begin{alertblock}{Ojo con esto}\small #1\end{alertblock}}
\newcommand{\apoyo}{{\normalfont\footnotesize\color{utpGris}~· material de apoyo}}
\newcommand{\marca}[1]{\tikz[baseline=(m.base)]{%
  \node[rounded corners=1.5pt,fill=utpGrisClaro,draw=utpGrisMedio,inner xsep=3pt,inner ysep=1.5pt,
        font=\footnotesize] (m) {#1};}}

\tikzset{
  caja/.style={draw=utpNegro, very thick, rounded corners=1.5pt, fill=white,
    align=center, minimum height=0.72cm, inner sep=5pt, font=\small},
  cajaazul/.style={caja, draw=utpGris, fill=white},
  cajaverde/.style={caja, draw=utpNegro, fill=utpGrisClaro},
  cajanaranja/.style={caja, draw=utpRojo!75, fill=utpRojo!7},
  flecha/.style={-{Stealth[length=2.3mm]}, very thick, draw=utpNegro},
  flechavuelta/.style={{Stealth[length=2.3mm]}-{Stealth[length=2.3mm]}, very thick, draw=utpNegro}
}

% ---------- Notación pata de gallo (crow's foot) ----------
% La rotación indica hacia dónde está la entidad: 0 derecha, 180 izquierda,
% 90 arriba, -90 abajo. La pata se abre hacia la entidad del lado "muchos".
\tikzset{
  relacion/.style={thick, draw=utpNegro},
  pics/uno/.style={code={%
    \draw[thick,utpNegro,line cap=round] (0,-1.5mm) -- (0,1.5mm);}},
  pics/muchos/.style={code={%
    \draw[thick,utpNegro,line cap=round,line join=round]
      (2.6mm,1.5mm) -- (0,0) -- (2.6mm,-1.5mm);
    \draw[thick,utpNegro,line cap=round] (0,0) -- (2.6mm,0);}},
}
\newcommand{\simbuno}{\tikz[baseline=-0.6ex]{\draw[relacion](0,0)--(7mm,0);\pic at (4.4mm,0){uno};}}
\newcommand{\simbmuchos}{\tikz[baseline=-0.6ex]{\draw[relacion](0,0)--(7mm,0);\pic at (4.4mm,0){muchos};}}


\title{Índices compuestos y consultas SELECT}
\subtitle{Índices compuestos, filtros en una tienda web\\ y estructura de una consulta completa}
\date{Semana 10 - Sesión 19}
\author{Rel Guzman}

% Espacio en blanco para completar
\newcommand{\bl}[1][1.9cm]{\underline{\hspace{#1}}}
% Nombre de tabla, siempre en naranja
\newcommand{\tbl}[1]{\textcolor{utpNaranja}{\texttt{#1}}}
% Palabra reservada de SQL dentro de un bloque de código
\newcommand{\kw}[1]{\textcolor{utpRojo}{\textbf{#1}}}
% Parte nueva de un bloque de codigo que crece paso a paso
\newcommand{\nuevo}[1]{{%
  \setlength{\fboxsep}{1.2pt}%
  \begingroup
  \renewcommand{\kw}[1]{\textcolor{utpNegro}{\textbf{##1}}}%
  \colorbox{yellow!45}{\textcolor{utpNegro}{#1}}%
  \endgroup}}
% Comentario dentro de un bloque de código SQL
\newcommand{\com}[1]{\textcolor{utpGris}{-- #1}}
% Marcas de clave primaria y clave foránea
\newcommand{\pk}{{\scriptsize\color{utpRojo}\textbf{PK}}}
\newcommand{\fk}{{\scriptsize\color{utpNaranja}\textbf{FK}}}
% Insignia de práctica
\newcommand{\ins}[1]{\tikz[baseline=-0.5ex]{\node[rounded corners=2pt,fill=utpRojo,text=white,
  font=\bfseries\scriptsize,inner xsep=5pt,inner ysep=2.5pt] {#1};}}
% Etiqueta oscura: nombre del gestor o de la forma de escribir algo
\newcommand{\etiqueta}[1]{\tikz[baseline=-0.5ex]{\node[rounded corners=2pt,fill=utpNegro,text=white,
  font=\bfseries\scriptsize,inner xsep=6pt,inner ysep=3pt] {#1};}}
% Celda resaltada en rojo claro
\newcommand{\mal}[1]{{\setlength{\fboxsep}{1.2pt}\colorbox{utpRojo!18}{#1}}}
% Barra invertida de las rutas de Windows
\newcommand{\bs}{\textbackslash{}}
% Bloque de código sobre fondo gris
\newsavebox{\sqlbuf}
\newenvironment{sqlcode}[1][.86\textwidth]
  {\begin{center}\begin{lrbox}{\sqlbuf}\begin{minipage}{#1}\ttfamily\footnotesize\vspace{6pt}\raggedright}
  {\vspace{6pt}\end{minipage}\end{lrbox}\colorbox{utpGrisClaro}{\usebox{\sqlbuf}}\end{center}}
% Sangría dentro de un bloque de código
\newcommand{\ind}{\hspace*{1.6em}}
% Anatomía de una instrucción: cada pieza del código con su pregunta al lado
\tikzset{
  pieza/.style={fill=utpGrisClaro,font=\ttfamily\footnotesize,anchor=west,inner sep=5pt,
    minimum height=0.62cm},
  para/.style={font=\small,anchor=west,text=utpNegro}
}
\newcommand{\piezax}{7.6}
\newcommand{\pieza}[3]{%
  \node[pieza] at (0,#1) {#2};
  \node[para] at (\piezax,#1) {\textcolor{utpRojo}{$\blacktriangleleft$}\ #3};}
% Diapositiva de sección: una palabra bien grande (\seccion[0]{...} la deja sin línea roja)
\newcommand{\seccion}[2][1]{%
\begin{frame}[plain]
  \begin{tikzpicture}[remember picture,overlay]
    \node[anchor=north east,inner sep=0pt]
      at ([xshift=-10mm,yshift=-8mm]current page.north east)
      {\includegraphics[width=43mm]{utp-logo.png}};
    \node[anchor=center,inner sep=0pt] at (current page.center)
      {\resizebox{100mm}{!}{\bfseries\color{utpNegro}#2}};
    \ifnum#1=1
    \node[anchor=north,inner sep=0pt] at ([yshift=-11mm]current page.center)
      {\color{utpRojo}\rule{100mm}{1.6pt}};
    \fi
  \end{tikzpicture}
\end{frame}}
% Cabecera de columna dentro de una tabla en monoespaciado
\newcommand{\cab}[1]{\normalfont\bfseries #1}
% Resultado de una consulta: \res{columnas}{cabecera}{filas}
\newcommand{\res}[3]{%
  \centering{\scriptsize\color{utpGris}resultado}\\[3pt]
  \footnotesize\ttfamily
  \begin{tabular}{#1}
  \toprule
  #2 \\
  \midrule
  #3
  \bottomrule
  \end{tabular}}
% Columna del modelo lógico: nombre y, en gris, su tipo
\newcommand{\col}[2]{\texttt{#1}~{\color{utpGris}#2}}
% Diagrama de la peluquería, igual en el modelo conceptual y en el lógico.
% \diagpelu[1 con verbos | 0 sin verbos]{cliente}{mascota}{cita}{pago}{servicio}{empleado}{encargado}
\newcommand{\diagpelu}[8][1]{%
  \ifnum#1=1 \def\entw{1.8cm}\else\def\entw{2.45cm}\fi
  \begin{tikzpicture}[font=\tiny]
    \tikzset{ent/.style={cajaverde,text width=\entw,align=left,font=\tiny,inner sep=4pt,
      minimum height=0pt},
      verbo/.style={font=\tiny\itshape,text=utpRojo,inner sep=1.5pt}}
    \node[ent] (cli) at (0,0) {#2};
    \node[ent] (mas) at (3.6,0) {#3};
    \node[ent] (cit) at (7.2,0) {#4};
    \node[ent] (pag) at (10.8,0) {#5};
    \node[ent] (ser) at (3.6,-3.0) {#6};
    \node[ent] (emp) at (10.8,-3.0) {#7};
    \node[ent,fill=white,draw=utpGris] (enc) at (0,-3.0) {#8};
    \coordinate (s) at ([xshift=-6mm]cit.south);
    \coordinate (e) at ([xshift=6mm]cit.south);
    % cliente 1 -- N mascota
    \draw[relacion] (cli.east) -- (mas.west);
    \pic at ($(cli.east)+(2.6mm,0)$) {uno};
    \pic at ($(mas.west)+(-2.6mm,0)$) {muchos};
    % mascota 1 -- N cita
    \draw[relacion] (mas.east) -- (cit.west);
    \pic at ($(mas.east)+(2.6mm,0)$) {uno};
    \pic at ($(cit.west)+(-2.6mm,0)$) {muchos};
    % cita 1 -- 1 pago
    \draw[relacion] (cit.east) -- (pag.west);
    \pic at ($(cit.east)+(2.6mm,0)$) {uno};
    \pic at ($(pag.west)+(-2.6mm,0)$) {uno};
    % servicio 1 -- N cita
    \draw[relacion] (ser.east) -| (s);
    \pic at ($(ser.east)+(2.6mm,0)$) {uno};
    \pic[rotate=90] at ($(s)+(0,-2.6mm)$) {muchos};
    % empleado 1 -- N cita
    \draw[relacion] (emp.west) -| (e);
    \pic at ($(emp.west)+(-2.6mm,0)$) {uno};
    \pic[rotate=90] at ($(e)+(0,-2.6mm)$) {muchos};
    \ifnum#1=1
      \path (cli.east) -- (mas.west) node[midway,above=1mm,verbo] {es dueño de};
      \path (mas.east) -- (cit.west) node[midway,above=1mm,verbo] {tiene};
      \path (cit.east) -- (pag.west) node[midway,above=1mm,verbo] {se paga con};
      \path (ser.east) -- (ser.east -| s) node[midway,above=0.5mm,verbo] {se presta en};
      \path (emp.west) -- (emp.west -| e) node[midway,above=0.5mm,verbo] {atiende};
    \fi
    \node[font=\tiny,text=utpGris,below=1.5mm of enc,text width=2.6cm,align=center]
      {no se relaciona: solo guarda el acceso};
  \end{tikzpicture}}

\begin{document}

% 1 ---------------------------------------------------------------------
{\setbeamertemplate{footline}{}\maketitle}

% 2 ---------------------------------------------------------------------
\begin{frame}[c]
\begin{block}{Logro de la sesión}
Al finalizar, podrás decidir cuándo conviene un \textbf{índice de varias columnas}, elegir el orden de sus columnas y leer la estructura de una \textbf{consulta \texttt{SELECT} completa}.
\end{block}
\end{frame}
% =======================================================================
% Revisión de la sesión anterior
% ---------------------------------------------------------------------
\begin{frame}{Revisión de la sesión anterior: en resumen}
\vspace{-0.2cm}
\begin{center}
\footnotesize
\renewcommand{\arraystretch}{1.2}
\begin{tabular}{>{\raggedright\bfseries}p{3.1cm}>{\raggedright\arraybackslash}p{4.9cm}>{\ttfamily\scriptsize\raggedright\arraybackslash}p{5.6cm}}
\toprule
& \bfseries En una frase & \normalfont\bfseries Ejemplo \\
\midrule
Índice & Una lista ordenada que dice dónde está cada fila. & CREATE INDEX idx\_cliente\_nombre ON cliente (nombre); \\
\texttt{IS NULL} & Pregunta si falta el dato. Nunca \texttt{= NULL}. & WHERE telefono IS NULL \\
\texttt{COALESCE} / \texttt{NULLIF} & Un repuesto para el \texttt{NULL}; convertir un valor en \texttt{NULL}. & COALESCE(telefono, 'sin teléfono') \\
Expresión & Un cálculo fila por fila, con su \texttt{AS}. & sueldo * 12 AS sueldo\_anual \\
\texttt{CASE WHEN} & Decide un valor según una condición. & CASE WHEN sueldo >= 3000 THEN 'alto' ELSE 'bajo' END \\
Agregación & Muchas filas, un solo valor. & COUNT(*), SUM, AVG, MIN, MAX \\
\texttt{GROUP BY} / \texttt{HAVING} & Un resultado por grupo; filtrar los grupos. & GROUP BY area HAVING COUNT(*) >= 2 \\
\bottomrule
\end{tabular}
\end{center}
\vspace{0.05cm}
\small Hoy repasamos dos ideas más a fondo: los \textbf{índices de varias columnas} y la \textbf{consulta \texttt{SELECT} completa}.
\end{frame}
% ---------------------------------------------------------------------
\begin{frame}{Revisión: un índice de \textbf{varias} columnas}
\begin{columns}[c,onlytextwidth]
\column{.54\textwidth}
\centering
\begin{tikzpicture}[y=0.56cm]
  \tikzset{
    celda/.style={draw=utpNegro,thick,font=\scriptsize\ttfamily,minimum height=0.48cm,
      inner sep=2.5pt,align=left},
    idx/.style={celda,fill=yellow!35,text width=2.05cm},
    dato/.style={celda,fill=utpGrisClaro,text width=2.35cm}
  }
  \node[idx] (i1) at (0,0) {Cómputo \ \ 45};
  \node[idx] (i2) at (0,-1) {Cómputo \ \ 3200};
  \node[idx] (i3) at (0,-2) {Hogar \ \ \ \ 90};
  \node[idx] (i4) at (0,-3) {Muebles \ \ 380};
  \node[idx] (i5) at (0,-4) {Muebles \ \ 650};
  \node[dato] (t1) at (4.3,0) {1 Laptop \ 3200};
  \node[dato] (t2) at (4.3,-1) {2 Mouse \ \ 45};
  \node[dato] (t3) at (4.3,-2) {3 Silla \ \ 380};
  \node[dato] (t4) at (4.3,-3) {4 Mesa \ \ \ 650};
  \node[dato] (t5) at (4.3,-4) {5 Lámpara \ 90};
  \draw[-{Stealth[length=1.8mm]},utpGris,thick] (i1.east)--(t2.west);
  \draw[-{Stealth[length=1.8mm]},utpGris,thick] (i2.east)--(t1.west);
  \draw[-{Stealth[length=1.8mm]},utpGris,thick] (i3.east)--(t5.west);
  \draw[-{Stealth[length=1.8mm]},utpRojo,very thick] (i4.east)--(t3.west);
  \draw[-{Stealth[length=1.8mm]},utpRojo,very thick] (i5.east)--(t4.west);
  \node[font=\scriptsize\bfseries,above=1mm of i1] {índice (categoria, precio)};
  \node[font=\scriptsize,above=1mm of t1] {\tbl{producto}};
\end{tikzpicture}
\column{.44\textwidth}
\small
\begin{itemize}\setlength\itemsep{7pt}
  \item Se ordena primero por \texttt{categoria}.
  \item Si dos tienen la \textbf{misma} categoría, se ordenan por \texttt{precio}.
  \item «Muebles de menos de 700» están \textbf{juntos} en la lista.
\end{itemize}
\end{columns}
\vspace{0.1cm}
\begin{sqlcode}[.72\textwidth]
\scriptsize
\kw{CREATE INDEX} idx\_producto\_categoria\_precio\\
\ind \kw{ON} producto \nuevo{(categoria, precio)};
\end{sqlcode}
\vspace{-0.15cm}
\analogia{Es la \textbf{guía telefónica}: ordenada por apellido y, dentro del mismo apellido, por nombre.}
\end{frame}
% ---------------------------------------------------------------------
\begin{frame}{Revisión: el orden de las columnas importa}
\small Con el índice \texttt{(categoria, precio)}, ¿qué consultas salen rápido?
\vspace{0.1cm}
\begin{center}
\small
\renewcommand{\arraystretch}{1.2}
\begin{tabular}{>{\ttfamily\footnotesize\raggedright\arraybackslash}p{7.0cm}>{\centering\arraybackslash}p{1.5cm}>{\raggedright\arraybackslash}p{4.5cm}}
\toprule
\normalfont\bfseries La consulta tiene\ldots & \bfseries ¿Lo usa? & \bfseries Por qué \\
\midrule
WHERE categoria = 'Muebles' & Sí & usa la \textbf{primera} columna \\
WHERE categoria = 'Muebles' AND precio < 700 & Sí & usa las \textbf{dos} columnas \\
WHERE categoria = 'Muebles' ORDER BY precio & Sí & ya vienen \textbf{ordenados} por precio \\
WHERE precio < 700 & \mal{No} & se \textbf{saltó} la primera columna \\
\bottomrule
\end{tabular}
\end{center}
\vspace{-0.1cm}
\concepto{La regla}{\small El índice se usa \textbf{de izquierda a derecha}. Pon primero la columna que se compara con \texttt{=}, y después la del rango o la del \texttt{ORDER BY}.}
\vspace{0.05cm}
\analogia{En la guía buscas rápido a los «Pérez» o a «Pérez, Ana». Pero no a todas las «Ana»: están repartidas por toda la guía.}
\end{frame}
% ---------------------------------------------------------------------
\begin{frame}{Revisión: una tienda web con filtros}
\begin{columns}[c,onlytextwidth]
\column{.47\textwidth}
\centering
\begin{tikzpicture}[font=\tiny]
  % ventana del navegador
  \draw[thick,utpNegro,rounded corners=2pt] (0,0) rectangle (6.6,-4.6);
  \fill[utpGrisClaro,rounded corners=2pt] (0.03,-0.03) rectangle (6.57,-0.45);
  \foreach \x in {0.2,0.38,0.56} \fill[utpGrisMedio] (\x,-0.24) circle (0.06);
  \node[draw=utpGrisMedio,fill=white,anchor=west,inner sep=1.5pt,text width=4.6cm,font=\tiny\ttfamily]
    at (0.85,-0.24) {localhost:3003};
  \node[anchor=west,font=\small\bfseries] at (0.2,-0.8) {Mi Tienda};
  % filtros
  \fill[utpGrisClaro] (0.2,-1.15) rectangle (2.2,-4.4);
  \node[anchor=west,font=\tiny\bfseries] at (0.3,-1.4) {Categoría};
  \node[draw,fill=yellow!35,anchor=west,text width=1.55cm,inner sep=2pt] at (0.3,-1.75)
    {Hogar \hfill$\blacktriangledown$};
  \node[anchor=west,font=\tiny\bfseries] at (0.3,-2.2) {Precio (S/)};
  \node[draw,fill=yellow!35,anchor=west,text width=0.55cm,inner sep=2pt] at (0.3,-2.55) {2000};
  \node[anchor=west] at (1.15,-2.55) {a};
  \node[draw,fill=yellow!35,anchor=west,text width=0.55cm,inner sep=2pt] at (1.4,-2.55) {2500};
  \node[anchor=west,font=\tiny\bfseries] at (0.3,-3.0) {Ordenar por};
  \node[draw,fill=yellow!35,anchor=west,text width=1.55cm,inner sep=2pt] at (0.3,-3.35)
    {menor precio \hfill$\blacktriangledown$};
  \node[fill=utpRojo,text=white,font=\tiny\bfseries,text width=1.55cm,align=center,inner sep=2.5pt,
    anchor=west] at (0.3,-3.95) {Filtrar};
  % productos
  \foreach \c in {0,1,2} \foreach \f in {0,1} {
    \draw[utpGrisMedio,rounded corners=1pt] (2.45+\c*1.4,-1.15-\f*1.65) rectangle ++(1.25,-1.45);
    \fill[utpGrisClaro] (2.55+\c*1.4,-1.25-\f*1.65) rectangle ++(1.05,-0.7);
    \node[anchor=west] at (2.5+\c*1.4,-2.15-\f*1.65) {Producto};
    \node[anchor=west,font=\tiny\bfseries] at (2.5+\c*1.4,-2.4-\f*1.65) {S/ 2000};
  }
\end{tikzpicture}
\column{.50\textwidth}
\small Cada filtro que elige el cliente se vuelve \textbf{una parte} de la consulta:
\begin{sqlcode}[\textwidth]
\scriptsize
\kw{SELECT} nombre, precio\\
\kw{FROM} producto\\
\kw{WHERE} \nuevo{categoria = 'Hogar'}\\
\ind \kw{AND} \nuevo{precio \kw{BETWEEN} 2000 \kw{AND} 2500}\\
\kw{ORDER BY} \nuevo{precio}\\
\kw{LIMIT} 10;
\end{sqlcode}
\vspace{-0.15cm}
\small Esta consulta se repite \textbf{cada vez} que alguien filtra. El índice que la acelera:
\begin{sqlcode}[\textwidth]
\scriptsize
\kw{CREATE INDEX} idx\_producto\_categoria\_precio\\
\ind \kw{ON} producto (\nuevo{categoria}, \nuevo{precio});
\end{sqlcode}
\end{columns}
\vspace{0.1cm}
\small\color{utpGris} Primero \texttt{categoria}, que va con \texttt{=}; después \texttt{precio}, que es un rango y además ordena.
\end{frame}
% ---------------------------------------------------------------------
\begin{frame}{Demostración: la tienda con un millón de productos}
\small La carpeta \texttt{tienda/} trae esa página. Con los botones \textbf{creas y borras} el índice, y comparas.
\vspace{0.2cm}
\begin{columns}[T,onlytextwidth]
\column{.48\textwidth}
\centering\etiqueta{Crear índice}
\begin{sqlcode}[\textwidth]
\scriptsize
\kw{CREATE INDEX} idx\_producto\_categoria\_precio\\
\ind \kw{ON} producto (categoria, precio);
\end{sqlcode}
\column{.48\textwidth}
\centering\etiqueta{Borrar índice}
\begin{sqlcode}[\textwidth]
\scriptsize
\kw{DROP INDEX} idx\_producto\_categoria\_precio;
\end{sqlcode}
\end{columns}
\vspace{0.15cm}
\begin{center}
\small
\renewcommand{\arraystretch}{1.35}
\begin{tabular}{>{\raggedright\arraybackslash}p{4.4cm}>{\raggedright\arraybackslash}p{4.0cm}>{\raggedright\arraybackslash}p{4.0cm}}
\toprule
& \bfseries Sin índice & \bfseries Con índice \\
\midrule
El filtro tarda & \bl[1.8cm] ms & \bl[1.8cm] ms \\
«Cómo buscó» dice & \texttt{SCAN producto} & \texttt{SEARCH \ldots{} USING INDEX} \\
\bottomrule
\end{tabular}
\end{center}
\vspace{0.1cm}
\small \texttt{SCAN} es leer la tabla \textbf{completa}, fila por fila. \texttt{SEARCH} es ir \textbf{directo} con el índice.
\vspace{0.2cm}

\small \textbf{Para probarlo}, en la carpeta \texttt{tienda/}: \texttt{npm install}, \texttt{npm run crear} y \texttt{npm run dev}. Luego abre \texttt{http://localhost:3003}.
\end{frame}
% ---------------------------------------------------------------------
\begin{frame}{\ins{Práctica 1} ¿Índice? ¿Sobre qué columnas?}
\small La tabla \tbl{pedido}\texttt{(id, cliente\_id, fecha, estado, total)} de una tienda en línea tiene \textbf{dos millones} de filas.
\vspace{0.15cm}
\begin{center}
\footnotesize
\renewcommand{\arraystretch}{1.25}
\begin{tabular}{>{\raggedright\arraybackslash}p{9.2cm}>{\raggedright\arraybackslash}p{3.8cm}}
\toprule
\bfseries Situación & \bfseries ¿Índice? ¿Sobre qué? \\
\midrule
1. La página «Mis pedidos»: \texttt{WHERE cliente\_id = ? ORDER BY fecha} & \bl[3.6cm] \\[4pt]
2. El almacén ve los pendientes: \texttt{WHERE estado = 'pendiente' AND fecha >= ?} & \bl[3.6cm] \\[4pt]
3. Abrir un pedido por su número: \texttt{WHERE id = ?} & \bl[3.6cm] \\[4pt]
4. Un reporte que se corre \textbf{una vez al año}: \texttt{WHERE total > 1000} & \bl[3.6cm] \\[4pt]
5. La tabla \tbl{metodo\_pago} tiene solo 3 filas: efectivo, tarjeta, Yape. & \bl[3.6cm] \\[4pt]
6. Ya existe \texttt{(cliente\_id, fecha)}. ¿Sirve para \texttt{WHERE fecha = ?} & \bl[3.6cm] \\
\bottomrule
\end{tabular}
\end{center}
\vspace{0.15cm}
\small \textbf{Para escribir:} el \texttt{CREATE INDEX} de las situaciones 1 y 2. \textbf{Pista:} el número de pedido es la clave primaria.
\end{frame}
% ---------------------------------------------------------------------
\begin{frame}{Revisión: todo junto, lo que puede llevar un \texttt{SELECT}}
\begin{center}
\begin{tikzpicture}[y=0.54cm]
  \renewcommand{\piezax}{7.4}
  \pieza{7}{\kw{SELECT} c.nombre, \kw{COUNT}(*) \kw{AS} pedidos}{qué columnas o cálculos se \textbf{muestran}}
  \pieza{6}{\kw{FROM} cliente c}{la tabla \textbf{principal}}
  \pieza{5}{\kw{JOIN} pedido p \kw{ON} c.id = p.cliente\_id}{qué otra tabla se \textbf{une}, y por dónde}
  \pieza{4}{\kw{WHERE} p.fecha >= '2026-01-01'}{qué \textbf{filas} entran}
  \pieza{3}{\kw{GROUP BY} c.nombre}{cómo se \textbf{agrupan}}
  \pieza{2}{\kw{HAVING} \kw{COUNT}(*) >= 2}{qué \textbf{grupos} se quedan}
  \pieza{1}{\kw{ORDER BY} pedidos \kw{DESC}}{en qué \textbf{orden} salen}
  \pieza{0}{\kw{LIMIT} 5;}{cuántas filas, \textbf{como máximo}}
\end{tikzpicture}
\end{center}
\vspace{-0.05cm}
\begin{itemize}\setlength\itemsep{3pt}
  \item \texttt{SELECT} indica qué mostrar. \texttt{FROM} se usa al consultar tablas. Lo demás se agrega si hace falta.
  \item Si se usan, se escriben \textbf{siempre en este orden}.
  \item \texttt{SELECT DISTINCT} quita las filas \textbf{repetidas} del resultado.
\end{itemize}
\vspace{0.05cm}
\small\color{utpGris} Se lee: «los 5 clientes con más pedidos este año, entre los que tienen 2 o más».
\end{frame}
% ---------------------------------------------------------------------
\begin{frame}{Cierre: lo de hoy en una tabla}
\begin{center}
\small
\renewcommand{\arraystretch}{1.45}
\begin{tabular}{>{\raggedright\bfseries}p{4.2cm}>{\raggedright\arraybackslash}p{8.6cm}}
\toprule
& \bfseries En una frase \\
\midrule
Índice de varias columnas & Ordena por la primera columna y, dentro de ella, por la siguiente. \\
Orden de las columnas & Importa el prefijo izquierdo del índice. \\
Tienda con filtros & Comparamos el plan y el tiempo de consulta con y sin índice. \\
\texttt{SELECT} completo & Cada cláusula tiene una función y un lugar al escribir la consulta. \\
\texttt{WHERE} / \texttt{HAVING} & Filtran filas y grupos, respectivamente. \\
\texttt{ORDER BY} / \texttt{LIMIT} & Ordenan el resultado y limitan la cantidad de filas. \\

\bottomrule
\end{tabular}
\end{center}
\vspace{0.25cm}
\small El índice se elige según las \textbf{consultas que hace la aplicación}.
\end{frame}
% ---------------------------------------------------------------------
\begin{frame}{Referencias}
\footnotesize
Coronel, C., Morris, S., \& Rob, P. (2011). \textit{Bases de datos: diseño, implementación y administración} (9.ª ed.). Cengage Learning.

\vspace{0.2cm}
Elmasri, R., \& Navathe, S. (2007). \textit{Fundamentos de sistemas de bases de datos} (5.ª ed.). Addison Wesley.

\vspace{0.2cm}
Silberschatz, A., Korth, H. F., \& Sudarshan, S. (2020). \textit{Database system concepts} (7th ed.). McGraw-Hill.

\vspace{0.2cm}
Microsoft. (s.\,f.). \textit{Transact-SQL reference}. \texttt{https://learn.microsoft.com/sql/t-sql/language-reference}

\vspace{0.2cm}
Microsoft. (s.\,f.). \textit{SQL Server Management Studio (SSMS)}. \texttt{https://learn.microsoft.com/sql/ssms}

\vspace{0.2cm}
Oracle. (s.\,f.). \textit{MySQL 8.0 reference manual: multiple-column indexes}. \texttt{https://dev.mysql.com/doc/refman/8.0/en/multiple-column-indexes.html}

\vspace{0.2cm}
SQLite. (s.\,f.). \textit{Query planning}. \texttt{https://www.sqlite.org/queryplanner.html}
\end{frame}
% ---------------------------------------------------------------------
% Cierre UTP
{
  \setbeamertemplate{background canvas}{%
    \includegraphics[width=\paperwidth,height=\paperheight]{gracias-final.pdf}%
  }
  \setbeamertemplate{footline}{}
  \begin{frame}[plain]
  \end{frame}
}

\end{document}
